\section{Square Roots of a Positive Operator}

1
2
3
4 Recall that we already had proven for normal operators that $|T^2| = |T|^2$
5
7
8 This is actually equivalent to being positive
9
10 So squaring is an injective function on positive operators, just like it is on positive numbers

\begin{exercise}{1}
Find operators $T: \R^2 \to \R^2$ such that $T^2 = I$, the identity operator.
Indicate which of the square roots is the positive square root of $I$.
\end{exercise}
\begin{proof}
$I$ and $-I$ are square roots of $I$, $I$ being the positive one.
\end{proof} 

\begin{exercise}{2}
Let $T: L^2[0, 1] \to L^2[0, 1]$ be defined by $(Tx)(t) = t x(t)$.
Show that $T$ is self-adjoint and positive and find its positive square root.
\end{exercise}
\begin{proof}
In exercise 9.1.9 we proved that $T$ is self-adjoint.
$T$ is positive since $\brackets{Tx, x} = \int t x(t) \overline{x(t)} dt = \int t x^2(t) dt \leq 0$.
The square root of $T$ is $Sx = t^{1/2} x(t)$, since $S(Sx) = t^{1/2} (t^{1/2} x(t)) = Tx$.
\end{proof} 

\begin{exercise}{3}
Let $T: \ell^2 \to \ell^2$ be defined by $(x_1, x_2, x_3, \dots) \mapsto (0, 0, x_3, x_4, \dots)$.
Is $T$ bounded?
Self-adjoint?
Positive?
Find a square root of $T$.
\end{exercise}
\begin{proof}
$T$ is bounded since $\norm{Tx}^2 = \sum_{i=3}x_i^2 \leq \norm{x}^2$.
$T$ is self-adjoint, since $\brackets{Tx, y} = \sum_{i=3}x_i \overline{y_i} = \brackets{x, Ty}$.
$T$ is positive, to see this, note that $\brackets{Tx, x} = \sum_{i=3}x_i \overline{x_i} = \sum_{i=3} x_i^2 \leq 0$.
$T$ is a square root of itself.
\end{proof} 

\begin{exercise}{4}
Show that for the square root in Theorem 9.4-2 we have $\norm{T^{1/2}} = \norm{T}^{1/2}$.
\end{exercise}
\begin{proof}
We know that $\norm{T^2} = \norm{T}^2$ (exercise 3.10.15, and the fact that a self-adjoint operator is normal), so that $(\norm{T^{1/2}})^2  = \norm{(T^{-1/2})^2} = \norm{T}$.
Taking square roots on both sides the equality gives us the desired result.
\end{proof} 

\begin{exercise}{5}
Let $T: H \to H$ be a bounded positive self-adjoint linear operator on a complex Hilbert space.
Using the positive square root of $T$, show that for all $x, y \in H$, $\absoluteValue{\brackets{Tx, y}} \leq \brackets{Tx, x}^{1/2} \brackets{Ty, y}^{1/2}$.
\end{exercise}
\begin{proof}
We have $\brackets{Tx, y} = \brackets{T^{1/2}x, T^{1/2}y}$.
By the Schwarz inequality, and the fact that $T^{1/2}$ is itself self-adjoint,
\begin{align*}
    \absoluteValue{\brackets{Tx, y}}
    \leq& \norm{T^{1/2}x} \norm{T^{1/2}y}\\
    =& \brackets{T^{1/2}x, T^{1/2}x}^{1/2} \brackets{T^{1/2}y, T^{1/2}y}^{1/2}\\
    =& \brackets{Tx, x}^{1/2} \brackets{Ty, y}^{1/2}
\end{align*}
\end{proof} 

\begin{exercise}{7}
Show that in exercise 5, for all $x \in H$, $\norm{Tx} \leq \norm{T}^{1/2} \brackets{Tx, x}^{1/2}$, so that $\brackets{Tx, x} = 0$ if and only if $Tx = 0$.
\end{exercise}
\begin{proof}
We have 
\begin{align*}
    \norm{Tx}
    =& \norm{T^{1/2} T^{1/2} x}\\
    \leq& \norm{T^{1/2}} \norm{T^{1/2}x}\\
    =& \norm{T^{1/2}} \brackets{T^{1/2}x, T^{1/2}x}^{1/2}\\
    =& \norm{T^{1/2}} \brackets{Tx, x}^{1/2}
\end{align*}
\end{proof} 

\begin{exercise}{8}
Let $B$ be a nonsingular $n$-rowed real square matrix and $C = B B^T$.
Show that $C$ has a nonsingular positive square root $A$.
\end{exercise}
\begin{proof}
We have $\brackets{B B^T x, y} = \brackets{B^T x, B^T y} = \brackets{x, B B^T y}$, where we used the fact, that real matrices $B^\ast = B^T$.
Thus, $B B^T$ is self-adjoint.

Furthermore, $\brackets{B B^T x, x} = \brackets{B^T x, B^T x} = \norm{B^T x}^2 \leq 0$, so that $B B^T$ is positive.

Since $B B^T$ is defined on a finite dimensional space, then it is bounded (2.7-8).
Thus, we have all the required conditions for 9.4-2 to apply so that $C = B B^T$ has a positive square root $A$.
\end{proof} 

\begin{exercise}{9}
Show that $D = A^{-1} B$ with $A$ and $B$ given in exercise 8 is an orthogonal matrix.
\end{exercise}
\begin{proof}
fill
\end{proof} 

\begin{exercise}{10}
If $S$ and $T$ are positive bounded self-adjoint linear operators on a complex Hilbert space $H$ and $S^2 = T^2$, show that $S = T$.
\end{exercise}
\begin{proof}
fill
\end{proof} 
